\documentclass[11pt,a4paper]{article}
\usepackage[margin=20mm,headheight=15pt]{geometry}
\usepackage{fontspec}
\setmainfont{TeX Gyre Pagella}
\setmonofont{DejaVu Sans Mono}[Scale=0.82]
\usepackage{microtype}
\usepackage{enumitem}
\usepackage{fvextra}
\usepackage{booktabs}
\usepackage{longtable}
\usepackage{fancyhdr}
\usepackage[hidelinks]{hyperref}
\setlength{\parindent}{0pt}
\setlength{\parskip}{6pt}
\setlength{\emergencystretch}{3em}
\hyphenpenalty=10000
\exhyphenpenalty=10000
\raggedright
\DefineVerbatimEnvironment{Code}{Verbatim}{fontsize=\small,
  breaklines=true,breakanywhere=true,breaksymbolleft={},breaksymbolright={},
  frame=single,framesep=5pt}
\setlist[enumerate]{leftmargin=7mm,itemsep=3pt,topsep=5pt,parsep=0pt}
\pagestyle{fancy}
\fancyhf{}
\fancyhead[L]{\small Computer Vision 1}
\fancyhead[R]{\small Exam 3 · Solutions}
\fancyfoot[C]{\small \thepage}
\hypersetup{pdftitle={Exam 3 — Solutions},pdfsubject={CV1 practice material}}
\begin{document}
{\LARGE\bfseries Exam 3 — Solutions}\par
{\large Synthesis and transfer}\par\bigskip
Answer key and worked explanations. Source aliases and convention warnings are listed in \texttt{COVERAGE.md} and \texttt{SOURCE\_NOTES.md} in the study pack.
\section*{Quick key}
\begin{longtable}{@{}r l@{\hspace{25mm}}r l@{}}
\toprule Question & Answer & Question & Answer \\ \midrule
01 & C & 26 & D \\
02 & C & 27 & A \\
03 & B & 28 & D \\
04 & B & 29 & A \\
05 & A & 30 & B \\
06 & B & 31 & 260.00 \\
07 & D & 32 & 0.26 \\
08 & A & 33 & 6.00 \\
09 & C & 34 & 0.25 \\
10 & D & 35 & 9 \\
11 & B & 36 & 10 \\
12 & D & 37 & 3 \\
13 & A & 38 & 5 \\
14 & B & 39 & 4.38 \\
15 & B & 40 & 14.62 \\
16 & C & 41 & 105 \\
17 & A & 42 & 2.33 \\
18 & D & 43 & 2.00 \\
19 & B & 44 & 0.50 \\
20 & A & 45 & 0.57 \\
21 & A & 46 & 0.45 \\
22 & C & 47 & 4448 \\
23 & C & 48 & 275 \\
24 & C & 49 & 0.36 \\
25 & D & 50 & 0.31 \\
\bottomrule\end{longtable}
\clearpage\section*{Worked solutions}
\begin{minipage}{\linewidth}
\textbf{Q01 — C}\par\smallskip
\textbf{C. n\_\allowbreak{}w =\allowbreak{} R.T @ n\_\allowbreak{}c; translation does not act on a direction vector.}\par
Normals between rigid coordinate frames transform by rotation only. R maps world directions to camera directions; its transpose maps back. A general non-rigid affine transform would require a different normal rule,\allowbreak{} which is not asked here.\par
{\small\itshape Topic: camera. Sources: HC1b pp.35–45; HC2a surface normals; prerequisite linear algebra.}

\end{minipage}
\par\vspace{12pt}
\begin{minipage}{\linewidth}
\textbf{Q02 — C}\par\smallskip
\textbf{C. High spatial frequencies can otherwise alias into false lower-frequency patterns.}\par
Subsampling reduces the representable frequency range. Filtering first suppresses frequencies that the coarse grid cannot represent; smoothing after aliasing does not undo it.\par
{\small\itshape Topic: filtering. Sources: HC2b pp.93–105; HC5b pp.38–39; Szeliski §3.4.}

\end{minipage}
\par\vspace{12pt}
\begin{minipage}{\linewidth}
\textbf{Q03 — B}\par\smallskip
\textbf{B. No; multiple albedo/\allowbreak{}shading combinations can produce the same image.}\par
For positive scale c,\allowbreak{} A'=\allowbreak{}cA and S'=\allowbreak{}S/\allowbreak{}c preserve I,\allowbreak{} and further spatial ambiguities exist. Priors or extra observations are required.\par
{\small\itshape Topic: reflectance. Sources: HC2a pp.4–6,\allowbreak{}21–23,\allowbreak{}47–48; Lab1.}

\end{minipage}
\par\vspace{12pt}
\begin{minipage}{\linewidth}
\textbf{Q04 — B}\par\smallskip
\textbf{B. The model assumes a linear relation between measured intensity and radiance; the response must be linearized or modelled.}\par
If recorded value is a nonlinear function of radiance,\allowbreak{} the original linear dot-product equations no longer describe the observations. Camera-response calibration or suitable raw data is needed.\par
{\small\itshape Topic: photometric. Sources: HC2a pp.19–20,\allowbreak{}61–62; Lab1.}

\end{minipage}
\par\vspace{12pt}
\begin{minipage}{\linewidth}
\textbf{Q05 — A}\par\smallskip
\textbf{A. The interaction of source spectrum,\allowbreak{} surface spectral reflectance and that channel's spectral sensitivity.}\par
A channel integrates wavelength-dependent contributions. Material colour alone does not fix measured RGB; illumination and sensor response matter.\par
{\small\itshape Topic: colour. Sources: HC1b pp.67–71,\allowbreak{}80–86; Lab1.}

\end{minipage}
\par\vspace{12pt}
\begin{minipage}{\linewidth}
\textbf{Q06 — B}\par\smallskip
\textbf{B. An additive constant in height.}\par
Derivatives determine relative height,\allowbreak{} not the reference height. Fixing a height such as z(0,\allowbreak{}0)=\allowbreak{}0 removes this additive ambiguity on a connected region.\par
{\small\itshape Topic: photometric. Sources: HC2a pp.41–43; Lab1; Forsyth–Ponce §2.2.}

\end{minipage}
\par\vspace{12pt}
\begin{minipage}{\linewidth}
\textbf{Q07 — D}\par\smallskip
\textbf{D. Interior outputs remain constant,\allowbreak{} but outputs near boundaries can be darker.}\par
Interior samples see a full constant patch and weights sum to 1. At a zero-padded boundary some weighted values are 0. Padding is part of the filter's actual behaviour.\par
{\small\itshape Topic: filtering. Sources: HC2b pp.47–87; Lab2.}

\end{minipage}
\par\vspace{12pt}
\begin{minipage}{\linewidth}
\textbf{Q08 — A}\par\smallskip
\textbf{A. It can fill sufficiently small gaps/\allowbreak{}holes relative to the element's shape,\allowbreak{} but need not fill a large cavity.}\par
Closing is scale- and shape-dependent. It is dilation followed by erosion,\allowbreak{} not general topological hole filling.\par
{\small\itshape Topic: morphology. Sources: HC2b pp.39–45; S2.}

\end{minipage}
\par\vspace{12pt}
\begin{minipage}{\linewidth}
\textbf{Q09 — C}\par\smallskip
\textbf{C. Its outermost pixel border is removed.}\par
For a border pixel,\allowbreak{} part of the tested support lies outside the image and encounters 0. Interior pixels whose entire support is inside survive.\par
{\small\itshape Topic: morphology. Sources: HC2b pp.27–38; S2.}

\end{minipage}
\par\vspace{12pt}
\begin{minipage}{\linewidth}
\textbf{Q10 — D}\par\smallskip
\textbf{D. Each coefficient is multiplied by a frequency-dependent phase factor of unit magnitude; its magnitude is unchanged.}\par
The translated spectrum is the original spectrum times a unit-magnitude complex exponential. This is exact for a circular shift; cropping/\allowbreak{}zero-filled shifts can also change magnitude through boundaries. A particular coefficient can have an unchanged phase when its factor is 1.\par
{\small\itshape Topic: fourier. Sources: HC2b pp.97–104; Szeliski §3.4.}

\end{minipage}
\par\vspace{12pt}
\begin{minipage}{\linewidth}
\textbf{Q11 — B}\par\smallskip
\textbf{B. They are complex conjugates,\allowbreak{} with indices interpreted modulo the image size.}\par
Real-valued signals have Hermitian symmetry F[-k]=\allowbreak{}conj(F[k]). Thus magnitudes are symmetric and phases have opposite signs modulo a full turn,\allowbreak{} where defined.\par
{\small\itshape Topic: fourier. Sources: HC2b pp.99–104; Szeliski §3.4.}

\end{minipage}
\par\vspace{12pt}
\begin{minipage}{\linewidth}
\textbf{Q12 — D}\par\smallskip
\textbf{D. Reflectance or illumination changes can also cause an intensity edge.}\par
Image derivatives measure appearance changes. Material boundaries,\allowbreak{} shadows,\allowbreak{} highlights and orientation changes can create edges without a depth break.\par
{\small\itshape Topic: features. Sources: HC3a; S3 introduction; HC2a.}

\end{minipage}
\par\vspace{12pt}
\begin{minipage}{\linewidth}
\textbf{Q13 — A}\par\smallskip
\textbf{A. To locate distinctive structures at their characteristic scales; DoG approximates a scale-normalized LoG response.}\par
Scale selection allows corresponding structures to be detected after resizing. DoG is an efficient approximation,\allowbreak{} not a guarantee of unlimited viewpoint invariance.\par
{\small\itshape Topic: matching. Sources: HC4a pp.36–42; S4; Szeliski §7.1.}

\end{minipage}
\par\vspace{12pt}
\begin{minipage}{\linewidth}
\textbf{Q14 — B}\par\smallskip
\textbf{B. Gradients remove the constant offset b,\allowbreak{} and descriptor normalization reduces the common gain a.}\par
The derivative of aI+b is a times the derivative of I. Orientation and normalized magnitudes are relatively stable,\allowbreak{} provided clipping,\allowbreak{} quantization and nonlinear response do not disrupt the assumption.\par
{\small\itshape Topic: matching. Sources: HC4a pp.41–44; Szeliski §7.1.2 pp.434–436.}

\end{minipage}
\par\vspace{12pt}
\begin{minipage}{\linewidth}
\textbf{Q15 — B}\par\smallskip
\textbf{B. It is geometrically degenerate and does not uniquely constrain a general 2D affine transform.}\par
All-inlier probability alone is not enough. A minimal sample must also have the geometry required to determine the model. For affine transforms,\allowbreak{} the [x,\allowbreak{}y,\allowbreak{}1] design rows must have rank 3.\par
{\small\itshape Topic: fitting. Sources: HC4a pp.17–22,\allowbreak{}45–55; Lab4.}

\end{minipage}
\par\vspace{12pt}
\begin{minipage}{\linewidth}
\textbf{Q16 — C}\par\smallskip
\textbf{C. The larger threshold can admit poor matches and outliers as inliers.}\par
A threshold controls what residual is accepted. Compare fit quality using meaningful units and tolerances,\allowbreak{} rather than maximizing a count under arbitrarily relaxed criteria.\par
{\small\itshape Topic: fitting. Sources: HC4a pp.46–48; HC3a pp.99–109.}

\end{minipage}
\par\vspace{12pt}
\begin{minipage}{\linewidth}
\textbf{Q17 — A}\par\smallskip
\textbf{A. The appearance-constancy assumption for a physical surface point is violated; apparent image motion need not equal physical surface motion.}\par
Optical flow describes apparent image change and relies on assumptions. Lighting changes can mimic motion even with stationary geometry. A moving shadow pattern may itself have a consistent image displacement; that is not the underlying surface's displacement.\par
{\small\itshape Topic: flow. Sources: HC3b pp.8–47,\allowbreak{}57–58; Lab3.}

\end{minipage}
\par\vspace{12pt}
\begin{minipage}{\linewidth}
\textbf{Q18 — D}\par\smallskip
\textbf{D. The object can disappear or become indistinct at coarse levels,\allowbreak{} so its motion is not initialized correctly.}\par
Pyramids reduce displacement but also reduce detail. Coarse estimation can miss thin or small structures,\allowbreak{} and later refinement may remain in the wrong solution basin.\par
{\small\itshape Topic: flow. Sources: HC3b pp.52–58.}

\end{minipage}
\par\vspace{12pt}
\begin{minipage}{\linewidth}
\textbf{Q19 — B}\par\smallskip
\textbf{B. Incorrect predictions can contaminate the model and cause tracking drift.}\par
Self-labelled updates can learn the occluder or background. Controlled updates and re-detection can help,\allowbreak{} but no method is universally immune.\par
{\small\itshape Topic: tracking. Sources: HC3b pp.89–106; Lab3.}

\end{minipage}
\par\vspace{12pt}
\begin{minipage}{\linewidth}
\textbf{Q20 — A}\par\smallskip
\textbf{A. Depth-dependent parallax; one global homography need not align the nonplanar scene.}\par
Translation makes image displacement depend on depth. Blending can hide seams but cannot generally correct incompatible geometric alignments or double images.\par
{\small\itshape Topic: geometry. Sources: HC4a pp.33–35,\allowbreak{}60–64; S4; Lab4.}

\end{minipage}
\par\vspace{12pt}
\begin{minipage}{\linewidth}
\textbf{Q21 — A}\par\smallskip
\textbf{A. T @ H}\par
Matrix composition applies the rightmost transformation first: source -> H -> reference -> T -> canvas. Transform direction and multiplication order must be explicit.\par
{\small\itshape Topic: geometry. Sources: HC4a pp.19,\allowbreak{}24–35; Lab4.}

\end{minipage}
\par\vspace{12pt}
\begin{minipage}{\linewidth}
\textbf{Q22 — C}\par\smallskip
\textbf{C. Unweighted 3-nearest-neighbour majority voting predicts class 0,\allowbreak{} although the closest example has class 1.}\par
Squared distances are [.02,\allowbreak{}.82,\allowbreak{}.82,\allowbreak{}16.82]. The three nearest labels are [1,\allowbreak{}0,\allowbreak{}0],\allowbreak{} in an irrelevant tie order for the two zero labels. Their sum is 1,\allowbreak{} so the comparison is false and prediction=\allowbreak{}0. k-NN uses labelled examples; it does not update cluster centres. With k=\allowbreak{}1 the prediction would instead be 1.\par
{\small\itshape Topic: classification. Sources: HC4b pp.37–53; Lab0.}

\end{minipage}
\par\vspace{12pt}
\begin{minipage}{\linewidth}
\textbf{Q23 — C}\par\smallskip
\textbf{C. Spatial terms encourage location coherence,\allowbreak{} and alpha sets their relative importance against colour differences.}\par
K-means uses the metric in the supplied feature space. Scaling positions changes distances and can trade colour similarity against spatial compactness; it does not guarantee semantics or connectivity.\par
{\small\itshape Topic: clustering. Sources: HC5a pp.91–118.}

\end{minipage}
\par\vspace{12pt}
\begin{minipage}{\linewidth}
\textbf{Q24 — C}\par\smallskip
\textbf{C. It permits margin violations while balancing fit against regularization.}\par
A hard-margin separator is infeasible for nonseparable labels. Soft margin penalizes violations and controls model complexity; it does not know which labels are wrong.\par
{\small\itshape Topic: classification. Sources: HC4b pp.69–78.}

\end{minipage}
\par\vspace{12pt}
\begin{minipage}{\linewidth}
\textbf{Q25 — D}\par\smallskip
\textbf{D. The probabilities are well calibrated to actual correctness frequencies.}\par
Softmax produces a normalized score distribution,\allowbreak{} but calibration is a separate empirical property. A high score alone is not proof that its stated frequency matches reality.\par
{\small\itshape Topic: metrics. Sources: HC5a p.5; HC5b p.78; prerequisite probability.}

\end{minipage}
\par\vspace{12pt}
\begin{minipage}{\linewidth}
\textbf{Q26 — D}\par\smallskip
\textbf{D. Prevent unintended accumulation of earlier batches' gradients before this batch's update.}\par
PyTorch gradient buffers accumulate unless cleared. Clearing gradients does not reset parameter values or disable gradient recording. Deliberate gradient accumulation is possible,\allowbreak{} but would need a different specified training schedule.\par
{\small\itshape Topic: optimization. Sources: Lab5 tutorial and training code.}

\end{minipage}
\par\vspace{12pt}
\begin{minipage}{\linewidth}
\textbf{Q27 — A}\par\smallskip
\textbf{A. A compatible shortcut,\allowbreak{} such as a stride-2 projection,\allowbreak{} so the tensors have matching shapes.}\par
Elementwise residual addition requires compatible dimensions. Identity shortcuts work only when the required shape matches; otherwise use an appropriate projection/\allowbreak{}downsampling path.\par
{\small\itshape Topic: architectures. Sources: HC6a pp.99–111.}

\end{minipage}
\par\vspace{12pt}
\begin{minipage}{\linewidth}
\textbf{Q28 — D}\par\smallskip
\textbf{D. IoU is unchanged,\allowbreak{} while an unnormalized squared coordinate-error loss generally scales by the square of that factor.}\par
Intersection and union areas both scale by the square factor,\allowbreak{} so their ratio stays fixed. Squared coordinate errors also scale by the square factor,\allowbreak{} without cancellation.\par
{\small\itshape Topic: detection. Sources: S5 Exercise 4; YOLOv4 §2.2.}

\end{minipage}
\par\vspace{12pt}
\begin{minipage}{\linewidth}
\textbf{Q29 — A}\par\smallskip
\textbf{A. The epipoles are at infinity in the horizontal baseline direction.}\par
The other camera centre projects to a point at infinity when the baseline is parallel to the image plane. Corresponding points lie on horizontal scanlines after rectification.\par
{\small\itshape Topic: stereo. Sources: HC6b pp.50–90.}

\end{minipage}
\par\vspace{12pt}
\begin{minipage}{\linewidth}
\textbf{Q30 — B}\par\smallskip
\textbf{B. SfM estimates camera poses and sparse structure; bundle adjustment refines poses and points jointly; MVS can densify using estimated cameras.}\par
SfM recovers geometry from correspondences and motion. Ordinary monocular reconstruction retains a global scale ambiguity without a metric reference. MVS and bundle adjustment have different roles.\par
{\small\itshape Topic: stereo. Sources: HC6b pp.108–149.}

\end{minipage}
\par\vspace{12pt}
\begin{minipage}{\linewidth}
\textbf{Q31 — 260.00}\par\smallskip
\textbf{Submit 260.00.} Unrounded reference: \texttt{260.0}.\par
After crop and resize: fy'=\allowbreak{}300 and cy'=\allowbreak{}(240-20)/\allowbreak{}2=\allowbreak{}110. v'=\allowbreak{}300*\allowbreak{}(1/\allowbreak{}2)+110=\allowbreak{}260. Extrinsics do not change.\par
{\small\itshape Topic: camera. Sources: HC1b pp.24–45; S1.}

\end{minipage}
\par\vspace{12pt}
\begin{minipage}{\linewidth}
\textbf{Q32 — 0.26}\par\smallskip
\textbf{Submit 0.26.} Unrounded reference: \texttt{0.2625}.\par
The per-pixel sums are [[10,\allowbreak{}10],\allowbreak{}[0,\allowbreak{}4]],\allowbreak{} with the zero denominator replaced by 1. Normalized red values are [[.6,\allowbreak{}.2],\allowbreak{}[0,\allowbreak{}.25]]. Their mean is 1.05/\allowbreak{}4=\allowbreak{}.2625; submit .26. Normalizing a spatially averaged RGB vector would be a different calculation.\par
{\small\itshape Topic: colour. Sources: Lab1 Section 1; Lab0 array axes.}

\end{minipage}
\par\vspace{12pt}
\begin{minipage}{\linewidth}
\textbf{Q33 — 6.00}\par\smallskip
\textbf{Submit 6.00.} Unrounded reference: \texttt{6.0}.\par
Path A: p=\allowbreak{}0 along y=\allowbreak{}0,\allowbreak{} then q=\allowbreak{}4 along x=\allowbreak{}2,\allowbreak{} giving z\_\allowbreak{}A=\allowbreak{}12. Path B: q=\allowbreak{}0 along x=\allowbreak{}0,\allowbreak{} then p=\allowbreak{}3 along y=\allowbreak{}3,\allowbreak{} giving z\_\allowbreak{}B=\allowbreak{}6. Difference=\allowbreak{}6. The field is not integrable because dp/\allowbreak{}dy=\allowbreak{}1 differs from dq/\allowbreak{}dx=\allowbreak{}2.\par
{\small\itshape Topic: photometric. Sources: HC2a pp.41–43; Lab1.}

\end{minipage}
\par\vspace{12pt}
\begin{minipage}{\linewidth}
\textbf{Q34 — 0.25}\par\smallskip
\textbf{Submit 0.25.} Unrounded reference: \texttt{0.25}.\par
n dot l =\allowbreak{} (-1/\allowbreak{}3)*\allowbreak{}.6+(2/\allowbreak{}3)*\allowbreak{}.8=\allowbreak{}1/\allowbreak{}3. It is positive,\allowbreak{} so I=\allowbreak{}.75/\allowbreak{}3=\allowbreak{}.25.\par
{\small\itshape Topic: reflectance. Sources: HC2a pp.12–20; S1; Lab1.}

\end{minipage}
\par\vspace{12pt}
\begin{minipage}{\linewidth}
\textbf{Q35 — 9}\par\smallskip
\textbf{Submit 9.} Unrounded reference: \texttt{9}.\par
Erosion leaves only the central pixel (4,\allowbreak{}4); the isolated corner pixel disappears. Dilation restores the central 3 by 3 block,\allowbreak{} giving 9 pixels. It does not restore the isolated noise.\par
{\small\itshape Topic: morphology. Sources: HC2b pp.39–42; S2.}

\end{minipage}
\par\vspace{12pt}
\begin{minipage}{\linewidth}
\textbf{Q36 — 10}\par\smallskip
\textbf{Submit 10.} Unrounded reference: \texttt{10}.\par
a+b=\allowbreak{}[0,\allowbreak{}10,\allowbreak{}10],\allowbreak{} whose median is 10. Both original medians are 0. The difference is 10,\allowbreak{} a concrete failure of additivity proving that median filtering is nonlinear.\par
{\small\itshape Topic: filtering. Sources: HC2b pp.47,\allowbreak{}89–91; Lab2.}

\end{minipage}
\par\vspace{12pt}
\begin{minipage}{\linewidth}
\textbf{Q37 — 3}\par\smallskip
\textbf{Submit 3.} Unrounded reference: \texttt{3}.\par
c[3]=\allowbreak{}1*\allowbreak{}h[3]+0*\allowbreak{}h[2]+2*\allowbreak{}h[1]+0*\allowbreak{}h[0]=\allowbreak{}1+2=\allowbreak{}3. Wrap-around is part of the defined operation.\par
{\small\itshape Topic: fourier. Sources: HC2b pp.101–103; Szeliski §3.4.}

\end{minipage}
\par\vspace{12pt}
\begin{minipage}{\linewidth}
\textbf{Q38 — 5}\par\smallskip
\textbf{Submit 5.} Unrounded reference: \texttt{5}.\par
The left group has one strong 10 and four connected weak 6s,\allowbreak{} so all 5 survive. The three 5s on the right have no connection to a strong seed and are discarded.\par
{\small\itshape Topic: features. Sources: HC3a pp.58–65; S2.}

\end{minipage}
\par\vspace{12pt}
\begin{minipage}{\linewidth}
\textbf{Q39 — 4.38}\par\smallskip
\textbf{Submit 4.38.} Unrounded reference: \texttt{4.381966011250105}.\par
M=\allowbreak{}[[6,\allowbreak{}1],\allowbreak{}[1,\allowbreak{}5]]. lambda\_\allowbreak{}min=\allowbreak{}(11-sqrt(5))/\allowbreak{}2=\allowbreak{}4.381966. Submit 4.38. The score tests the least constrained direction.\par
{\small\itshape Topic: features. Sources: S3; Lab3 Sections 2–3.}

\end{minipage}
\par\vspace{12pt}
\begin{minipage}{\linewidth}
\textbf{Q40 — 14.62}\par\smallskip
\textbf{Submit 14.62.} Unrounded reference: \texttt{14.615915982243466}.\par
Scale coarse displacement by 8: [14,\allowbreak{}-4]. Add the residual: [14.25,\allowbreak{}-3.25]. Magnitude=\allowbreak{}sqrt(14.25\textasciicircum{}2+3.25\textasciicircum{}2)=\allowbreak{}about 14.62 pixels. Resizing the array alone is not enough.\par
{\small\itshape Topic: flow. Sources: HC3b pp.48–54; Lab3.}

\end{minipage}
\par\vspace{12pt}
\begin{minipage}{\linewidth}
\textbf{Q41 — 105}\par\smallskip
\textbf{Submit 105.} Unrounded reference: \texttt{105}.\par
w\textasciicircum{}3=\allowbreak{}.064. log(.001)/\allowbreak{}log(.936)=\allowbreak{}104.4417. Ceiling gives 105. A homography would require a different sample size and a larger budget.\par
{\small\itshape Topic: fitting. Sources: HC3a pp.104–105; Lab4.}

\end{minipage}
\par\vspace{12pt}
\begin{minipage}{\linewidth}
\textbf{Q42 — 2.33}\par\smallskip
\textbf{Submit 2.33.} Unrounded reference: \texttt{2.3333333333333335}.\par
The first pair gives c=\allowbreak{}1,\allowbreak{}f=\allowbreak{}2. The next pairs imply a=\allowbreak{}3g+2,\allowbreak{}d=\allowbreak{}2g,\allowbreak{}b=\allowbreak{}h,\allowbreak{}e=\allowbreak{}4h+2. The last pair gives h=\allowbreak{}0,\allowbreak{}g=\allowbreak{}1. Thus H=\allowbreak{}[[5,\allowbreak{}0,\allowbreak{}1],\allowbreak{}[2,\allowbreak{}2,\allowbreak{}2],\allowbreak{}[1,\allowbreak{}0,\allowbreak{}1]]. At (.5,\allowbreak{}.5),\allowbreak{} x'=\allowbreak{}3.5/\allowbreak{}1.5=\allowbreak{}7/\allowbreak{}3; submit 2.33. Averaging target corner positions would not implement a projective map.\par
{\small\itshape Topic: geometry. Sources: HC4a pp.47–55; S4.}

\end{minipage}
\par\vspace{12pt}
\begin{minipage}{\linewidth}
\textbf{Q43 — 2.00}\par\smallskip
\textbf{Submit 2.00.} Unrounded reference: \texttt{2.0}.\par
Ordinary ratios are [2/\allowbreak{}3,\allowbreak{}3/\allowbreak{}4,\allowbreak{}1/\allowbreak{}2,\allowbreak{}5/\allowbreak{}6,\allowbreak{}1/\allowbreak{}2]. Accepted nearest distances are [2,\allowbreak{}1,\allowbreak{}.5],\allowbreak{} whose maximum is 2. Comparing D1/\allowbreak{}D2 directly with .7 would use the wrong threshold.\par
{\small\itshape Topic: matching. Sources: HC4a pp.43–44; S4; Lab4.}

\end{minipage}
\par\vspace{12pt}
\begin{minipage}{\linewidth}
\textbf{Q44 — 0.50}\par\smallskip
\textbf{Submit 0.50.} Unrounded reference: \texttt{0.5}.\par
Assignments produce counts [1,\allowbreak{}2,\allowbreak{}3]. After normalization by 6,\allowbreak{} the final bin is 3/\allowbreak{}6=\allowbreak{}.5. The fixed vocabulary gives the same histogram length for every image.\par
{\small\itshape Topic: classification. Sources: HC4b pp.23–28; S4.}

\end{minipage}
\par\vspace{12pt}
\begin{minipage}{\linewidth}
\textbf{Q45 — 0.57}\par\smallskip
\textbf{Submit 0.57.} Unrounded reference: \texttt{0.5714285714285714}.\par
The first four predictions are positive. TP=\allowbreak{}2,\allowbreak{} FP=\allowbreak{}2,\allowbreak{} FN=\allowbreak{}1. F1=\allowbreak{}4/\allowbreak{}(4+2+1)=\allowbreak{}4/\allowbreak{}7=\allowbreak{}.5714286; submit .57. The score exactly at .6 is positive because the rule is inclusive.\par
{\small\itshape Topic: metrics. Sources: HC4b pp.80–86; S4.}

\end{minipage}
\par\vspace{12pt}
\begin{minipage}{\linewidth}
\textbf{Q46 — 0.45}\par\smallskip
\textbf{Submit 0.45.} Unrounded reference: \texttt{0.45}.\par
Cumulative precisions are [0,\allowbreak{}1/\allowbreak{}2,\allowbreak{}1/\allowbreak{}3,\allowbreak{}2/\allowbreak{}4,\allowbreak{}3/\allowbreak{}5],\allowbreak{} recalls [0,\allowbreak{}1/\allowbreak{}4,\allowbreak{}1/\allowbreak{}4,\allowbreak{}2/\allowbreak{}4,\allowbreak{}3/\allowbreak{}4]. The future-maximum precision at each TP is .6. Three increments of .25 contribute .25*\allowbreak{}.6 each,\allowbreak{} so AP=\allowbreak{}.45. The missing final quarter of recall contributes zero; raw uninterpolated AP would instead be .40.\par
{\small\itshape Topic: metrics. Sources: HC5a pp.68–78; S4.}

\end{minipage}
\par\vspace{12pt}
\begin{minipage}{\linewidth}
\textbf{Q47 — 4448}\par\smallskip
\textbf{Submit 4448.} Unrounded reference: \texttt{4448}.\par
Counts are (64+1)*\allowbreak{}16=\allowbreak{}1040,\allowbreak{} (3*\allowbreak{}3*\allowbreak{}16+1)*\allowbreak{}16=\allowbreak{}2320,\allowbreak{} and (16+1)*\allowbreak{}64=\allowbreak{}1088. Total=\allowbreak{}4448. The identity shortcut has no parameters.\par
{\small\itshape Topic: architectures. Sources: HC5b pp.70–73; HC6a pp.107–111; S5.}

\end{minipage}
\par\vspace{12pt}
\begin{minipage}{\linewidth}
\textbf{Q48 — 275}\par\smallskip
\textbf{Submit 275.} Unrounded reference: \texttt{275}.\par
Each filter has 3*\allowbreak{}3*\allowbreak{}3*\allowbreak{}2=\allowbreak{}54 weights. Five filters and five biases give (54+1)*\allowbreak{}5=\allowbreak{}275. Do not count frames as independent input channels for this 3D convolution.\par
{\small\itshape Topic: cnn. Sources: S5 Exercise 5.}

\end{minipage}
\par\vspace{12pt}
\begin{minipage}{\linewidth}
\textbf{Q49 — 0.36}\par\smallskip
\textbf{Submit 0.36.} Unrounded reference: \texttt{0.36}.\par
Box confidence=\allowbreak{}.75*\allowbreak{}.6=\allowbreak{}.45. Class-specific score=\allowbreak{}.8*\allowbreak{}.45=\allowbreak{}.36. The convention is explicitly given because different YOLO versions use different outputs and targets.\par
{\small\itshape Topic: detection. Sources: HC6a pp.66–72; SOURCE\_\allowbreak{}NOTES.}

\end{minipage}
\par\vspace{12pt}
\begin{minipage}{\linewidth}
\textbf{Q50 — 0.31}\par\smallskip
\textbf{Submit 0.31.} Unrounded reference: \texttt{0.3125}.\par
Sensitivity magnitude is 180/\allowbreak{}144=\allowbreak{}1.25 metres per pixel. Multiply by .25: .3125 metres,\allowbreak{} so submit .31. Depth becomes increasingly sensitive to disparity errors at small disparity/\allowbreak{}large depth.\par
{\small\itshape Topic: stereo. Sources: HC6b pp.84–90; Szeliski §12.1; prerequisite calculus.}

\end{minipage}
\par\vspace{12pt}

\end{document}
